\documentclass[11pt]{article}

% ---------- Encoding & fonts ----------
\usepackage[T1]{fontenc}
\usepackage[utf8]{inputenc}

% ---------- Geometry ----------
\usepackage[left=1.85in,right=0.8in,top=0.6in,bottom=0.7in]{geometry}

% ---------- Packages ----------
\usepackage{titlesec}
\usepackage{enumitem}
\usepackage{fancyhdr}
\usepackage[hidelinks]{hyperref}
\usepackage[expansion=false]{microtype}

% ---------- Editable metadata ----------
\newcommand{\cvname}{Santiago Rodríguez Vázquez}
\newcommand{\cvdate}{%
  \ifcase\month\or Enero\or Febrero\or Marzo\or Abril\or Mayo\or Junio\or
    Julio\or Agosto\or Septiembre\or Octubre\or Noviembre\or Diciembre\fi {/}\number\year}

\newcommand{\pagenumcentered}{%
  \hspace*{\dimexpr 0.5\paperwidth - \oddsidemargin - 1in\relax}%
  \makebox[0pt][c]{\footnotesize\thepage}%
}

% ---------- Marginal section headers ----------
\titleformat{\section}[leftmargin]
  {\normalfont\bfseries\filright}
  {}{0pt}{}
\titlespacing{\section}{1.35in}{1.3em}{0.6em}

% ---------- Page style ----------
\pagestyle{fancy}
\fancyhf{}
\renewcommand{\headrulewidth}{0pt}
%\fancyfoot[L]{\pagenumcentered}

\fancypagestyle{firstpage}{%
  \fancyhf{}
  \renewcommand{\headrulewidth}{0pt}
%  \fancyfoot[L]{\pagenumcentered}
}

% ---------- Spacing helpers ----------
\setlength{\parindent}{0pt}
\newcommand{\entrydate}[1]{\hfill{\small{#1}}}
\newlist{cvitems}{itemize}{1}
\setlist[cvitems]{leftmargin=1.2em,label=\textendash,itemsep=1.5pt,topsep=2pt,parsep=0pt}

\newcommand{\grp}[1]{\smallskip{\itshape #1}\par\nopagebreak\vspace{1pt}}

\begin{document}
\thispagestyle{firstpage}

% ================= HEADER =================
{\LARGE\bfseries {\cvname}}

\href{mailto:santiago@rodriguez-vazquez.com.ar}{santiago@rodriguez-vazquez.com.ar} | \href{https://sanarovaz.github.io}{sanarovaz.github.io} | \href{https://www.linkedin.com/in/srodriguezvazquez/?locale=en-US}{linkedin.com/srodriguezvazquez}
\\
+54 3442 516481

% ================= PERFIL =================
\section{Perfil}
Economista con sólida formación cuantitativa (Maestría en curso) y experiencia práctica en análisis de datos, gestión de bases financieras e investigación aplicada. Combina herramientas econométricas y de programación (R, Python, SQL, Stata) con experiencia previa en el sector corporativo. Perfil analítico con varios años de experiencia docente.

% ================= FORMACIÓN =================
\section{Formación}
\textbf{Universidad de San Andrés}.\entrydate{2026--Presente}\\
\textit{Maestría en Economía.}
\begin{cvitems}
  \item Cursado de economía aplicada, comercio internacional, macroeconomía y econometría aplicada / organización industrial.
\end{cvitems}

\smallskip
\textbf{Universidad de Buenos Aires, Facultad de Ciencias Económicas}.\entrydate{2019--2025}\\
\textit{Licenciatura en Economía.}
\begin{cvitems}
  \item Promedio: 8.03/10. Materias cursadas del ciclo de Actuario (Estadística Actuarial, Análisis Numérico).
  \item Tesis de grado sobre los determinantes del financiamiento de startups espaciales, calificada con 10/10.
\end{cvitems}

% ================= EXPERIENCIA PROFESIONAL =================
\section{Experiencia\\profesional}
\textbf{Veritran}.\entrydate{2022--2023}\\
\textit{Analista Financiero Jr.}
\begin{cvitems}
  \item Gestión de operaciones de cuentas a pagar para subsidiarias en Latinoamérica, Estados Unidos y Europa, garantizando un procesamiento preciso y oportuno.
  \item Mantenimiento y actualización de bases de datos y registros financieros en NetSuite, dando soporte al reporting y al cierre mensual.
\end{cvitems}

% ================= INVESTIGACIÓN / ANÁLISIS DE DATOS =================
\section{Investigación y\\análisis de datos}
\textbf{Instituto Interdisciplinario de Economía Política (IIEP)}.\entrydate{2024--2026}\\
\textit{Asistente de Investigación.}
\begin{cvitems}
  \item Construcción de una base de datos financiera de 156 multilatinas para un estudio financiado por el BID sobre expansión regional (\href{https://publications.iadb.org/es/multilatinas-en-movimiento-caracteristicas-motivaciones-y-lineamientos-de-politica}{Blyde et al., 2025}).
  \item Aplicación de modelos estadísticos de redes (ERGM, QAP) en R para analizar datos a nivel firma sobre transferencia de conocimiento entre clústers tecnológicos.
  \item Colaboración en investigación y redacción sobre la historia del sector espacial argentino.
\end{cvitems}

\smallskip
\textbf{Proyecto independiente}, con Paulo Pascuini.\entrydate{2025--Presente}
\begin{cvitems}
  \item Análisis de patrones de financiamiento en 170 startups de nanosatélites y 410 rondas de inversión, estudiando el rol de la integración vertical y el perfil del inversor.
\end{cvitems}

% ================= DOCENCIA =================
\section{Docencia}
\textbf{Ayudante de Cátedra}, Universidad de San Andrés y Universidad de Buenos Aires.\entrydate{2022--Presente}
\begin{cvitems}
  \item Siete cuatrimestres como ayudante en cinco materias de economía (Microeconomía I y II, Economía Internacional, Economía Financiera, Desarrollo Económico), preparando material, dando clases y corrigiendo exámenes.
\end{cvitems}

% ================= HERRAMIENTAS =================
\section{Herramientas}
    R, Python. Manejo de SQL, Stata, Tableau Public. Algunos proyectos en \href{https://www.github.com/sanarovaz/}{GitHub}.
\section{Idiomas}
Español (nativo). Inglés (avanzado): Cambridge CAE, puntaje 198, 2017.

\end{document}